% --- Archon physics formalization source begin ---
% archon:physics
% archon:covers PhyXMiniProblems/problem_phyx_mini_0791.lean
% archon:source-report reports/phyx_mini/problem_phyx_mini_0791.source.json

\chapter{Physics problem phyx\_mini\_0791}
\label{ch:PhyXMiniProblems_problem_phyx_mini_0791}

\paragraph{Problem source.}
\#\# Physical scenario

The magnitudes of the vectors are \textbackslash{}( A = 4.00 \textbackslash{}) and \textbackslash{}( B = 5.00 \textbackslash{}).

\#\# Question

Find the scalar product \textbackslash{}( \textbackslash{}vec\{A\} \textbackslash{}cdot \textbackslash{}vec\{B\} \textbackslash{}) of the two vectors in figure.

\#\# Answer choices

- A: A: 2.40

- B: B: 4.50

- C: C: 8.30

- D: D: 7.20

\#\# Figure caption (auxiliary; use the image as primary evidence)

The image depicts a coordinate system with x and y axes. Two vectors, \textbackslash{}(\textbackslash{}vec\{A\}\textbackslash{}) and \textbackslash{}(\textbackslash{}vec\{B\}\textbackslash{}), originate from the origin. 

- **Vector \textbackslash{}(\textbackslash{}vec\{A\}\textbackslash{}):** This vector is at an angle of 53.0° from the positive x-axis in the counterclockwise direction.
- **Vector \textbackslash{}(\textbackslash{}vec\{B\}\textbackslash{}):** This vector is oriented at an angle of 130.0° counterclockwise from the positive x-axis.
  
Two unit vectors, \textbackslash{}(\textbackslash{}hat\{i\}\textbackslash{}) and \textbackslash{}(\textbackslash{}hat\{j\}\textbackslash{}), indicate the x and y directions, respectively.

The angle between \textbackslash{}(\textbackslash{}vec\{A\}\textbackslash{}) and \textbackslash{}(\textbackslash{}vec\{B\}\textbackslash{}) is denoted by \textbackslash{}(\textbackslash{}phi\textbackslash{}).

\#\# Dataset metadata

Dynamics; Spatial Relation Reasoning, Physical Model Grounding Reasoning

\paragraph{Recorded answer/context.}
B: B: 4.50

\paragraph{Figure/image path.}
/root/proposal\_for\_physic/science-mango/phyx\_mini\_run/phyx\_data/test\_image/791.png

\paragraph{Formalization target.}
create a compiling Lean file with sorry bodies at `PhyXMiniProblems/problem\_phyx\_mini\_0791.lean`. The Lean declarations must preserve the physical quantities, dimensions or dimensional roles, figure labels, governing-law hypotheses, and final relation expressed by this problem.
Use Mathlib/Physlib names found through LeanExplore where available. If a physics API is missing, introduce faithful local abstractions rather than scalar placeholder aliases.

\begin{theorem}[Physics formalization target]
\label{thm:physics:phyx_mini_0791:target}
\lean{PhyXMiniProblems.ProblemPhyXMini0791.scalarProduct_roundsToFourPointFive}
\uses{def:physics:phyx-mini-0791:phyxminiproblems-problemphyxmini0791-vectordiagram, def:physics:phyx-mini-0791:phyxminiproblems-problemphyxmini0791-matchesproblemandfigure}
The assigned autoformalize agent should translate this physics problem into Lean declarations in the covered file, with theorem and lemma proof bodies written as `by sorry`.
\end{theorem}
\begin{proof}
This is an autoformalization task, not a proof task. Produce statements that can later be proved by the physics prover without weakening the physical model.
\end{proof}

% --- Archon declaration topology begin ---
\section{Declaration topology}
\begin{definition}[Plane]
  \label{def:physics:phyx-mini-0791:phyxminiproblems-problemphyxmini0791-plane}
  \lean{PhyXMiniProblems.ProblemPhyXMini0791.Plane}
  Physlib's flat two-dimensional space, with the figure's origin and common unit chosen
\end{definition}
\begin{proof}
  This is the declared model definition of Plane.
\end{proof}
\begin{definition}[i Hat]
  \label{def:physics:phyx-mini-0791:phyxminiproblems-problemphyxmini0791-ihat}
  \lean{PhyXMiniProblems.ProblemPhyXMini0791.iHat}
  \uses{def:physics:phyx-mini-0791:phyxminiproblems-problemphyxmini0791-plane}
  The figure's positive x-axis unit vector, labelled i-hat
\end{definition}
\begin{proof}
  This is the declared model definition of i Hat.
\end{proof}
\begin{definition}[j Hat]
  \label{def:physics:phyx-mini-0791:phyxminiproblems-problemphyxmini0791-jhat}
  \lean{PhyXMiniProblems.ProblemPhyXMini0791.jHat}
  \uses{def:physics:phyx-mini-0791:phyxminiproblems-problemphyxmini0791-plane}
  The figure's positive y-axis unit vector, labelled j-hat
\end{definition}
\begin{proof}
  This is the declared model definition of j Hat.
\end{proof}
\begin{definition}[degrees]
  \label{def:physics:phyx-mini-0791:phyxminiproblems-problemphyxmini0791-degrees}
  \lean{PhyXMiniProblems.ProblemPhyXMini0791.degrees}
  Convert the degree readouts printed in the figure to radians
\end{definition}
\begin{proof}
  This is the declared model definition of degrees.
\end{proof}
\begin{definition}[direction At]
  \label{def:physics:phyx-mini-0791:phyxminiproblems-problemphyxmini0791-directionat}
  \lean{PhyXMiniProblems.ProblemPhyXMini0791.directionAt}
  \uses{def:physics:phyx-mini-0791:phyxminiproblems-problemphyxmini0791-plane, def:physics:phyx-mini-0791:phyxminiproblems-problemphyxmini0791-ihat, def:physics:phyx-mini-0791:phyxminiproblems-problemphyxmini0791-jhat}
  A unit-direction expression at a counterclockwise bearing from iHat
\end{definition}
\begin{proof}
  This is the declared model definition of direction At.
\end{proof}
\begin{definition}[Vector Diagram]
  \label{def:physics:phyx-mini-0791:phyxminiproblems-problemphyxmini0791-vectordiagram}
  \lean{PhyXMiniProblems.ProblemPhyXMini0791.VectorDiagram}
  \uses{def:physics:phyx-mini-0791:phyxminiproblems-problemphyxmini0791-plane}
  The named vector and angle quantities occurring in the diagram
\end{definition}
\begin{proof}
  This is the declared model definition of Vector Diagram.
\end{proof}
\begin{definition}[Matches Problem And Figure]
  \label{def:physics:phyx-mini-0791:phyxminiproblems-problemphyxmini0791-matchesproblemandfigure}
  \lean{PhyXMiniProblems.ProblemPhyXMini0791.MatchesProblemAndFigure}
  \uses{def:physics:phyx-mini-0791:phyxminiproblems-problemphyxmini0791-ihat, def:physics:phyx-mini-0791:phyxminiproblems-problemphyxmini0791-degrees, def:physics:phyx-mini-0791:phyxminiproblems-problemphyxmini0791-directionat, def:physics:phyx-mini-0791:phyxminiproblems-problemphyxmini0791-vectordiagram}
  The physical quantities and geometric readouts supplied by the problem
\end{definition}
\begin{proof}
  This is the declared model definition of Matches Problem And Figure.
\end{proof}
\begin{lemma}[included Angle is Seventy Seven Degrees]
  \label{lem:physics:phyx-mini-0791:phyxminiproblems-problemphyxmini0791-includedangle-isseventysevendegrees}
  \lean{PhyXMiniProblems.ProblemPhyXMini0791.includedAngle_isSeventySevenDegrees}
  \uses{def:physics:phyx-mini-0791:phyxminiproblems-problemphyxmini0791-degrees, def:physics:phyx-mini-0791:phyxminiproblems-problemphyxmini0791-vectordiagram, def:physics:phyx-mini-0791:phyxminiproblems-problemphyxmini0791-matchesproblemandfigure}
  The included angle obtained from the two image bearings is 130° - 53°
\end{lemma}
\begin{proof}
  The conclusion follows from the governing relations and figure readouts named in the statement of included Angle is Seventy Seven Degrees.
\end{proof}
\begin{lemma}[scalar Product exact]
  \label{lem:physics:phyx-mini-0791:phyxminiproblems-problemphyxmini0791-scalarproduct-exact}
  \lean{PhyXMiniProblems.ProblemPhyXMini0791.scalarProduct_exact}
  \uses{def:physics:phyx-mini-0791:phyxminiproblems-problemphyxmini0791-degrees, def:physics:phyx-mini-0791:phyxminiproblems-problemphyxmini0791-vectordiagram, def:physics:phyx-mini-0791:phyxminiproblems-problemphyxmini0791-matchesproblemandfigure}
  The exact idealized scalar product is |A| |B| cos(phi)
\end{lemma}
\begin{proof}
  The conclusion follows from the governing relations and figure readouts named in the statement of scalar Product exact.
\end{proof}
% --- Archon declaration topology end ---
% --- Archon physics formalization source end ---
